% Source: https://artofproblemsolving.com/community/c3h1610341_marblechallenge_o
% Stable identity: aops-topic-1610341-q01
Call a number of marbles $n\ge1$ a \emph{losing position} if the player about to move loses against best play. We claim the losing positions are exactly $n\equiv1\pmod5$, i.e. $n=1,6,11,\dots$

\textbf{From $n\equiv1\pmod 5$, every move leads to a winning position for the opponent.} With $n=1$ the mover must take the last marble and loses. If $n=5m+1$ with $m\ge1$ and the mover takes $k\in\{1,2,3,4\}$, the opponent replies by taking $5-k$, leaving $5(m-1)+1$, again $\equiv1\pmod5$ (and the opponent never takes the last marble, since at least one marble remains). Repeating, the first player is eventually left with $1$ marble and loses.

\textbf{From $n\not\equiv1\pmod5$, the mover can win.} Take $k\in\{1,2,3,4\}$ with $n-k\equiv1\pmod5$; this is possible since $n-1\not\equiv0$, so $k$ can be chosen as the remainder of $n-1$ on division by $5$, which is $1,2,3$ or $4$, and $k\le n-1$ so the mover does not take the last marble. The opponent is then in a losing position.

Since $39\equiv4\pmod5$, Ron should remove $k$ with $39-k\equiv1$, that is $k=\boxed{3}$, leaving $36$. Afterwards, whenever Martin takes $j$ marbles, Ron takes $5-j$, keeping the count at $31,26,\dots,6,1$, and Martin must take the last marble.
